% TEMPLATE-MILC-v3.tex - plantilla maestra UNICA de las guias MILC v3.
%
% La usan las cuatro familias de guias, cada una con su driver:
%   · Grados 6-11          -> scripts/build-guias-g11.py
%   · Bebras               -> scripts/build-guias-bebras.py
%   · Semillero            -> scripts/build-guias-semillero.py
%   · Territorio interior  -> scripts/build-guias-territorio-interior.py
%
% Antes habia tres copias casi identicas (~400 lineas iguales, 6 distintas):
% cada mejora habia que aplicarla tres veces, y en la practica no se hacia
% (el motor de recursos vivia solo en dos de ellas). Lo que varia entre
% familias esta parametrizado en 6 placeholders que cada driver llena
% (se nombran aqui SIN su sintaxis real: el builder reemplaza por texto plano,
% tambien dentro de los comentarios, y un valor multilinea rompe el preambulo):
%
%   HEADER_IZQ        encabezado izquierdo de pagina
%   HEADER_DER        encabezado derecho de pagina
%   PORTADA_ETIQUETA  rotulo sobre el numero grande (GRADO/MOMENTO/MODULO)
%   PORTADA_SERIE     linea de serie de la portada
%   PORTADA_ID        identificador de la guia en la portada
%   PORTADA_CIRCULO   el circulito de grado (°), o vacio si no aplica
%   PDF_TITULO        titulo en texto plano (sin \textbf ni comillas TeX) para
%                     los metadatos del PDF (/Title); lo produce lib_xelatex.texto_pdf
%
% Composicion (auditoria 2026-09-03): las bandas de estacion piden espacio
% (\Needspace*) y prohiben el salto justo despues (\nopagebreak), asi no quedan
% huerfanas al pie; las cajas largas (softbox, drawbox, infoband, puentebox) son
% `breakable`, asi una caja de media pagina ya no arrastra todo a la siguiente
% dejando la anterior al 40 %. Todo texto sobre mostaza o verde va en milcNegro
% (blanco sobre #E5B400 da 1,93:1 y sobre #58A923 2,95:1; ninguno pasa WCAG AA).
% El resto de claves las documenta content/guias/_SCHEMA.md. El script valida
% que no queden placeholders sin reemplazar antes de escribir el archivo final.

% Etiquetado PDF (latex-lab): /Lang, estructura y texto alternativo de las
% figuras, para lectores de pantalla. Debe ir ANTES de \documentclass.
\DocumentMetadata{lang=es-CO,tagging=on}
\documentclass[11pt,letterpaper]{article}
% headheight=24pt: la cabecera derecha lleva dos lineas (identidad de la guia
% y estacion actual). headsep se acorta en la misma medida para que el bloque
% de texto quede exactamente donde estaba con headheight=12pt.
\usepackage[margin=2.0cm,headheight=24pt,headsep=13pt]{geometry}
\usepackage[table]{xcolor}
\usepackage{fontspec}
\usepackage[spanish]{babel}
\usepackage{tikz}
% Librerías TikZ para los diagramas de las guías (bloque `recursos:`):
%  · babel  → babel en español vuelve activos caracteres como «>», y TikZ los usa
%             en sus opciones (>=stealth, ->). Sin esta librería, cualquier
%             diagrama con flechas revienta con «\language@active@arg> has an
%             extra }». Debe ir ANTES de que se dibuje nada.
%  · positioning        → sintaxis «below left=2pt and 3pt».
%  · decorations.pathreplacing → llaves ({brace}) para señalar rangos.
\usetikzlibrary{babel,positioning,decorations.pathreplacing}
\usepackage{graphicx}
% Circuitos: el trabajo de electrónica se hace hoy con micro:bit y sus sensores
% integrados (MakeCode, sin cableado), así que no se necesitan esquemáticos.
% Si algún día entran montajes con protoboard, basta descomentar esta línea:
% los diagramas .tex/.tikz declarados en `recursos:` ya se incrustan solos.
% \usepackage{circuitikz}
\usepackage[most]{tcolorbox}
\usepackage{tabularx}
\usepackage{array}
\usepackage{fancyhdr}
\usepackage{titlesec}
\usepackage[document]{ragged2e}  % bandera a la derecha en todo el cuerpo
\usepackage{enumitem}
\usepackage{microtype}
\usepackage{etoolbox}
\usepackage{needspace}
% hyperref al final del preambulo: metadatos del PDF (titulo, autor, idioma,
% familia), marcadores por estacion (\pdfbookmark) y enlaces sin recuadro.
\usepackage[hidelinks,
  pdftitle={Tablas limpias — campos, registros y una bitácora de lo que corregiste},
  pdfauthor={PhD. Álvaro Cárdenas Orozco},
  pdfsubject={Tecnología e Informática · Grado 8 · Guía 3 · MILC v3},
  pdflang={es-CO},
  bookmarksopen=true,bookmarksnumbered=false]{hyperref}

% Helvetica Neue no trae la flecha U+2192, asi que cada «→» escrito literalmente
% en el YAML se compilaba a NADA, en silencio (462 flechas en 61 guias: rutas del
% tipo «paleta → tipografia → lockup» salian rotas). Mapeamos el caracter a su
% version matematica, que si tiene glifo. Asi el autor puede seguir escribiendo
% «→» comodamente en el YAML.
\usepackage{newunicodechar}
\newunicodechar{→}{\ensuremath{\rightarrow}}
% Mismo problema con los simbolos de marca y vinetas: el «✓» de «una palomita ✓»
% salia como cuadrito vacio en el PDF de todas las guias que lo usaban.
\usepackage{pifont}
\newunicodechar{✓}{\ding{51}}
\newunicodechar{✗}{\ding{55}}
\newunicodechar{✔}{\ding{52}}
\newunicodechar{•}{\textbullet}
\newunicodechar{≥}{\ensuremath{\geq}}
\newunicodechar{≤}{\ensuremath{\leq}}

\setmainfont{Helvetica Neue}
\setsansfont{Helvetica Neue}
\setlength{\parindent}{0pt}
\setlength{\parskip}{6pt}
% Sin particion de palabras: en 6.º-7.º una palabra cortada por guion al final
% de linea («Comuni-cacion») frena la lectura mucho mas que una linea corta.
\hyphenpenalty=10000
\exhyphenpenalty=10000
\pretolerance=2000
\tolerance=3000
\setlength{\emergencystretch}{3em}
\linespread{1.10}
\renewcommand{\arraystretch}{1.25}
\setlist{leftmargin=5mm,itemsep=1mm,topsep=1mm}
\newcolumntype{Y}{>{\RaggedRight\arraybackslash}X}

\definecolor{milcMagenta}{HTML}{E6007E}
\definecolor{milcVino}{HTML}{5A0038}
\definecolor{milcTurquesa}{HTML}{0093A5}
\definecolor{milcVerde}{HTML}{58A923}
\definecolor{milcMostaza}{HTML}{E5B400}
\definecolor{milcOcre}{HTML}{B86600}
\definecolor{milcNegro}{HTML}{241816}
\definecolor{milcGris}{HTML}{F7F7F4}
\definecolor{milcLinea}{HTML}{E8E2DD}
\definecolor{milcGradePrimary}{HTML}{FF6600}
\definecolor{milcGradeDark}{HTML}{9A3E00}
\definecolor{milcGradeSoft}{HTML}{FFE4D0}
\definecolor{milcGradeAccent}{HTML}{CFE7FF}
\definecolor{milcGradeText}{HTML}{FFFFFF}
\color{milcNegro}

\pagestyle{fancy}
\fancyhf{}
% Cabecera de dos lineas: identidad de la guia arriba y, a la derecha y en
% gris, la estacion en curso (\markright desde \milcestacion). Antes de la
% primera estacion la segunda linea va vacia; el \strut mantiene la altura
% para que la linea de arriba no baile entre paginas.
\lhead{\small Tecnología e Informática · Grado 8\\[-1pt]{\footnotesize\strut}}
\rhead{\small Guía 3 · MILC v3\\[-1pt]{\footnotesize\strut\color{milcNegro!65}\rightmark}}
\cfoot{\small\thepage}
\renewcommand{\headrulewidth}{0pt}

% Sobre mostaza (#E5B400) y verde (#58A923) el texto va en milcNegro; sobre el
% resto de la paleta, en blanco. Se usa dentro de fonttitle/fontupper, donde un
% \color posterior manda sobre coltitle.
\newcommand{\milctintasobre}[1]{%
  \ifboolexpr{test{\ifstrequal{#1}{milcMostaza}} or test{\ifstrequal{#1}{milcVerde}}}%
    {\color{milcNegro}}{\color{white}}}

\newtcolorbox{titlebox}[1]{
  enhanced,colback=#1,colframe=#1,arc=7mm,boxrule=0pt,
  left=7mm,right=7mm,top=3mm,bottom=3mm,
  before skip=8pt,after skip=8pt,
  fontupper=\sffamily\bfseries\Large\color{white}
}

\newtcolorbox{softbox}[3]{
  enhanced,breakable,pad at break=3mm,
  colback=#2,colframe=#1,borderline west={4pt}{0pt}{#1},
  arc=2mm,boxrule=.4pt,left=4mm,right=4mm,top=3mm,bottom=3mm,
  before skip=6pt,after skip=8pt,title={#3},coltitle=white,
  colbacktitle=#1,fonttitle=\sffamily\bfseries\milctintasobre{#1},fontupper=\RaggedRight,
  attach boxed title to top left={xshift=3mm,yshift=-2mm},
  boxed title style={arc=2mm, boxrule=0pt}
}

\newtcolorbox{drawbox}[2]{
  enhanced,breakable,pad at break=4mm,
  colback=white,colframe=#1,arc=4mm,boxrule=1pt,
  left=5mm,right=5mm,top=4mm,bottom=4mm,before skip=8pt,after skip=8pt,
  title={#2},coltitle=white,colbacktitle=#1,fonttitle=\sffamily\bfseries\milctintasobre{#1},
  fontupper=\RaggedRight,
  attach boxed title to top left={xshift=4mm,yshift=-2mm},
  boxed title style={arc=3mm, boxrule=0pt}
}

\newtcolorbox{infoband}[1]{
  enhanced,breakable,pad at break=2mm,colback=#1!8,colframe=#1!50,arc=2mm,boxrule=.5pt,
  left=4mm,right=4mm,top=2mm,bottom=2mm,before skip=4pt,after skip=4pt,
  fontupper=\small\RaggedRight
}

% Puente narrativo entre secciones (contrato editorial Conéctate).
% Da continuidad: "Acabas de X. Ahora vas a Y porque Z."
% Visualmente: banda izquierda en color del grado, texto en itálica pequeña.
\newtcolorbox{puentebox}{
  enhanced,breakable,pad at break=2mm,colback=milcGradeSoft,colframe=milcGradeDark,
  borderline west={3pt}{0pt}{milcGradeDark},
  arc=0mm,boxrule=0pt,left=5mm,right=4mm,top=2.5mm,bottom=2.5mm,
  before skip=4pt,after skip=8pt,
  fontupper=\itshape\small\color{milcNegro}\RaggedRight
}
% La flecha va en modo matematico a proposito: Helvetica Neue NO tiene el glifo
% U+2192, asi que \textrightarrow se compilaba a NADA («Missing character: There
% is no → in font Helvetica Neue») y el puente salia sin flecha. En modo
% matematico la flecha viene de la fuente matematica y si se dibuja.
\newcommand{\puente}[1]{%
  \begin{puentebox}%
    \textcolor{milcGradeDark}{$\rightarrow$}\hspace{2mm}#1%
  \end{puentebox}%
}

\newcommand{\linead}[1]{\noindent\color{milcLinea}\rule{#1}{0.5pt}\color{milcNegro}}
\newcommand{\respuesta}[1]{\par\vspace{2mm}\foreach \n in {1,...,#1}{\noindent\color{milcLinea}\rule{\linewidth}{0.5pt}\par\vspace{5mm}}\color{milcNegro}}
\newcommand{\checkbox}{\textcolor{milcGradeDark}{\fbox{\phantom{x}}}\hspace{5pt}}

% ─── Listas aireadas para las guías socioemocionales ────────────────────
% Componer las verificaciones como «\checkbox texto\\» deja el texto que salta
% de línea debajo del cuadrito y apelmaza el bloque. Estos entornos alinean la
% continuación y airean los ítems. Los usa Territorio Interior; cualquier
% familia puede hacerlo.
\newlist{checklista}{itemize}{1}
\setlist[checklista]{
  label=\textcolor{milcGradeDark}{\fbox{\phantom{x}}},
  leftmargin=9mm, labelsep=3.2mm, itemsep=2.4mm,
  topsep=2.5mm, parsep=0pt, partopsep=0pt, before=\RaggedRight
}
\newlist{pasoslista}{enumerate}{1}
\setlist[pasoslista]{
  label=\textcolor{milcGradeDark}{\sffamily\bfseries\arabic*.},
  leftmargin=8mm, labelsep=3mm, itemsep=2.8mm,
  topsep=1mm, parsep=0pt, partopsep=0pt, before=\RaggedRight
}

\newcommand{\phasecard}[4]{
\begin{tcolorbox}[enhanced,colback=#3,colframe=#2,arc=3mm,boxrule=.5pt,height=3.05cm,valign=center,halign=center,left=1.5mm,right=1.5mm,top=2mm,bottom=2mm]
{\sffamily\bfseries\fontsize{22}{24}\selectfont\textcolor{#2}{#1}}\par
{\sffamily\bfseries\small #4}
\end{tcolorbox}
}

% ── Piezas del rediseno 2026 (legibilidad 6.º-11.º) ───────────────────
% Banda de estacion: reemplaza el titlebox macizo. Titulo a la izquierda,
% dato practico (tiempo/modalidad) a la derecha, en una sola linea.
% Nunca huerfana: pide 9 lineas libres (o salta de pagina rellenando con
% \vfill) y prohibe el salto justo despues de la banda. Nueve y no seis:
% la caja partible que sigue necesita ~80pt (titulo adosado, relleno y dos
% lineas) para arrancar en la misma pagina; con menos, tcolorbox la manda
% entera a la siguiente y la banda se queda sola. El penalty 10000 va
% en `after`, ANTES del espacio posterior: un `after skip` (aunque sea 0pt)
% es un \vskip tras la caja, y ese glue es un punto de corte legal que un
% \nopagebreak escrito despues de \end{tcolorbox} ya no cierra. Cada banda
% deja marcador PDF y actualiza la cabecera derecha.
\newcounter{milcestacion}
\newcommand{\milcestacion}[2]{%
\Needspace*{9\baselineskip}%
\stepcounter{milcestacion}%
\pdfbookmark[1]{#2}{milcestacion\themilcestacion}%
\markright{#2}%
\begin{tcolorbox}[enhanced,colback=#1,colframe=#1,arc=2mm,boxrule=0pt,
  left=5mm,right=5mm,top=2.6mm,bottom=2.6mm,before skip=11pt,
  after={\par\nopagebreak\vskip7pt}]
{\sffamily\bfseries\fontsize{15}{18}\selectfont\milctintasobre{#1}#2}
\end{tcolorbox}}

% Tarjeta de paso numerada: sustituye las tablas «Pilar / Aplicacion», que
% eran formato de consulta docente, no de estudiante. El numero va en un
% circulo sobre el borde para que la secuencia se lea de un vistazo.
\newcommand{\milcpaso}[3]{%
\Needspace{4\baselineskip}%
\begin{tcolorbox}[enhanced,colback=white,colframe=milcLinea,arc=1.5mm,boxrule=.9pt,
  left=14mm,right=4mm,top=3mm,bottom=3mm,before skip=4pt,after skip=4pt,
  fontupper=\RaggedRight,
  overlay={\node[anchor=north west,circle,fill=#2,text=white,inner sep=0pt,
    minimum size=7.4mm,font=\sffamily\bfseries\small]
    at ([xshift=3.6mm,yshift=-3.2mm]frame.north west) {#1};}]
#3
\end{tcolorbox}}

% Casilla marcable de tamano util con lapiz (la anterior era \fbox{\phantom{x}},
% de alto variable segun el texto que la seguia).
\newcommand{\milcmarca}{\framebox[3.8mm]{\rule{0pt}{3.4mm}}}

% Fila marcable de lista suelta (checks de la estacion 1).
\newcommand{\milccheckrow}[1]{%
\par\vspace{1.8mm}\noindent
\makebox[6.4mm][l]{\raisebox{-.55mm}{\textcolor{milcNegro}{\milcmarca}}}%
\parbox[t]{\dimexpr\linewidth-6.4mm\relax}{\RaggedRight #1}\par}

% Celda marcable dentro de la rubrica (tabla de autoevaluacion). Misma
% sangria francesa, pero pensada para vivir dentro de una columna X.
\newcommand{\milcopcion}[1]{%
\makebox[5.2mm][l]{\raisebox{-.55mm}{\textcolor{milcNegro}{\milcmarca}}}%
\parbox[t]{\dimexpr\linewidth-5.2mm\relax}{\RaggedRight #1}}

% ── Recursos visuales de la guía ──────────────────────────────────────
% Declarados en el bloque `recursos:` del YAML y emitidos por el builder.
% \guiaFigura   → imagen rasterizada o PDF (foto, carta celeste, montaje).
% \guiaDiagrama → fuente TikZ/circuitikz (.tex) incrustada como vector:
%                 es la vía para circuitos y esquemáticos editables.
% [#1] = texto alternativo (alt) · #2 = ruta relativa al .tex · #3 = pie
% El alt llega al PDF etiquetado (/Alt) y lo lee el lector de pantalla.
\newcommand{\guiaFigura}[3][]{%
\begin{center}
\includegraphics[width=0.88\linewidth,height=0.38\textheight,keepaspectratio,alt={#1}]{#2}\\[1.5mm]
{\footnotesize\itshape\textcolor{milcNegro!75}{#3}}
\end{center}
}

\newcommand{\guiaDiagrama}[2]{%
\begin{center}
\input{#1}\\[1.5mm]
{\footnotesize\itshape\textcolor{milcNegro!75}{#2}}
\end{center}
}

\begin{document}
\thispagestyle{empty}

% ── Portada ───────────────────────────────────────────────────────────
% Antes: un rectangulo de color a pagina completa. Se veia bien en pantalla,
% pero en la fotocopiadora del colegio se comia el toner y salia gris sucio.
% Ahora la banda ocupa el tercio superior y el resto va en blanco.
\begin{tikzpicture}[remember picture,overlay]
  \path[fill=milcGradePrimary]
    (current page.north west) rectangle ([yshift=-10.6cm]current page.north east);

  \node[anchor=north west,text=milcGradeText] at ([xshift=2.0cm,yshift=-1.55cm]current page.north west)
    {\sffamily\bfseries\fontsize{12}{14}\selectfont GRADO};
  \node[anchor=north west,text=milcGradeText,opacity=.85] at ([xshift=2.0cm,yshift=-2.35cm]current page.north west)
    {\sffamily\bfseries\fontsize{8.5}{10}\selectfont SERIE GUÍAS · TECNOLOGÍA E INFORMÁTICA};
  \node[anchor=north east,text=milcGradeText,opacity=.85,align=right] at ([xshift=-2.0cm,yshift=-1.55cm]current page.north east)
    {\sffamily\bfseries\fontsize{9}{12}\selectfont I.E. Sor María Juliana\\Cartago, Valle del Cauca};

  \node[anchor=west,text=milcGradeText] (gradonum) at ([xshift=1.85cm,yshift=-7.1cm]current page.north west)
    {\sffamily\bfseries\fontsize{112}{112}\selectfont 8};
  % El circulito de grado (°) solo aplica a las guias de grado («11°»); en
  % Bebras y Semillero el numero grande es un momento/modulo, no un grado.
  % Se posiciona relativo al nodo `gradonum`, no en coordenadas absolutas.
  \draw[milcGradeText,line width=1.6pt]
    ([xshift=2.0mm,yshift=-4.5mm]gradonum.north east) circle (.15cm);

  \node[anchor=west,text=milcGradeText,text width=11.4cm,align=left] at ([xshift=1.5cm,yshift=0.5cm]gradonum.east)
    {
     {\sffamily\bfseries\fontsize{9}{11}\selectfont\MakeUppercase{Período 1 · Análisis de datos con phronesis}}\\[3mm]
     {\sffamily\bfseries\fontsize{21}{25}\selectfont\setlength{\baselineskip}{27pt}Tablas limpias\\[4pt]y su bitácora}\\[3.5mm]
     {\sffamily\bfseries\fontsize{15}{18}\selectfont Guía 3-8-TIC}
    };
\end{tikzpicture}

\vspace*{9.4cm}
\vfill

{\sffamily\bfseries\fontsize{8}{10}\selectfont\textcolor{milcGradeDark}{LO QUE TE LLEVAS HOY}}

\begin{tcolorbox}[enhanced,colback=white,colframe=milcNegro,arc=2mm,boxrule=1.6pt,
  left=5mm,right=5mm,top=4mm,bottom=4mm,before skip=3pt,after skip=12pt,
  fontupper=\RaggedRight\fontsize{11}{15.5}\selectfont]
Una tabla de Excel con al menos 20 registros y 5 campos sobre un tema real del colegio, con formato de tabla aplicado. Y una bitácora que documente al menos cinco errores corregidos, con captura de antes y de después de cada uno.
\end{tcolorbox}

\begin{tcolorbox}[enhanced,colback=milcGris,colframe=milcGris,arc=2mm,boxrule=0pt,
  left=5mm,right=5mm,top=3.5mm,bottom=3.5mm,after skip=12pt]
{\sffamily\bfseries\fontsize{8}{10}\selectfont\textcolor{milcNegro!55}{ESTA GUÍA ES DE}}\\[3mm]
\begin{tabularx}{\linewidth}{X p{3.2cm} p{3.2cm}}
\linead{\linewidth} & \linead{3.0cm} & \linead{3.0cm}\\[-1mm]
{\fontsize{8.5}{10}\selectfont\textcolor{milcNegro!55}{Nombre completo}} &
{\fontsize{8.5}{10}\selectfont\textcolor{milcNegro!55}{Grupo}} &
{\fontsize{8.5}{10}\selectfont\textcolor{milcNegro!55}{Fecha}}\\
\end{tabularx}
\end{tcolorbox}

\vfill

\noindent\textcolor{milcLinea}{\rule{\linewidth}{1pt}}\\[2mm]
{\sffamily\bfseries\fontsize{9.5}{12}\selectfont Autor: Dr. Álvaro Cárdenas Orozco}\\[1mm]
{\sffamily\fontsize{8.5}{11}\selectfont\textcolor{milcNegro!55}{Metodología MILC · apertura ancestral · pensamiento computacional · triángulo Dussel–estoicismo–Floridi}}

\newpage

\section*{Apertura: Tablas limpias --- campos, registros y una bitácora de lo que corregiste}

\begin{softbox}{milcGradeDark}{milcGradeSoft}{Saber ancestral}
En los pueblos indígenas del Cauca, las normas que rigen a todos se llaman \textbf{mandatos}. No las escribe una autoridad sola. Salen de un congreso al que llega la gente de los resguardos a discutir durante días. El XV Congreso del Consejo Regional Indígena del Cauca se reunió en Río Blanco, Sotará, del 25 al 30 de junio de 2017 (CRIC, 2017). De allí salieron mandatos concretos, numerados y escritos. Y lo más importante viene después: en el siguiente congreso, la organización rinde cuentas en público de cuáles se cumplieron y cuáles no. El registro se revisa, se corrige y se vuelve a mirar. Un mandato no vale por estar escrito; vale porque quienes lo hicieron vuelven a mirarlo. La \textbf{cara de exclusión}: un mandato es un acto de gobierno de un pueblo que se disputa con el Estado y con actores armados. El propio CRIC reporta que muchos quedan sin cumplir. Nosotros hacemos hoy algo mucho más pequeño que aprende de eso: una tabla con sus registros y una bitácora de lo que corregimos, para que cualquiera pueda volver a mirarla.\par\smallskip{\footnotesize\textcolor{milcNegro!70}{\textbf{Fuente.} Consejo Regional Indígena del Cauca. (2017). \emph{Avanza el XV congreso regional del CRIC en el resguardo de Rioblanco Sotará del pueblo indígena yanacona}.}}
\end{softbox}

\begin{softbox}{milcTurquesa}{milcGris}{Saber contemporáneo}
Una \textbf{tabla de datos} en Excel tiene dos partes. Cada columna es un \textbf{campo}: una característica que se mide, como nombre, edad, fecha o nota. Cada fila es un \textbf{registro}: una persona o un evento concreto, con todos sus campos. Para que la tabla sirva, tiene que estar limpia: sin espacios extra, sin «María» y «MARIA» como si fueran dos personas, sin fechas mal escritas, sin filas repetidas y sin celdas vacías. Hay una regla vieja en el oficio de los datos: \emph{basura entra, basura sale}. Si la tabla está sucia, ningún cálculo posterior sirve. Y hay una costumbre que separa al aficionado del profesional: la \textbf{bitácora de limpieza}. Cada corrección se anota: qué error había, qué hiciste, cuándo y por qué. Sin bitácora, nadie puede volver a mirar tu trabajo. Ni tú.
\end{softbox}

\begin{softbox}{milcMostaza}{milcGris}{Pregunta-puente}
\textit{El CRIC no solo escribe los mandatos: en el siguiente congreso rinde cuentas de cuáles cumplió. Cuando limpies una tabla, ¿vas a poder decir qué corregiste, cuándo y por qué? ¿O solo vas a tener una tabla que «se ve bien»?}
\end{softbox}

\begin{softbox}{milcVerde}{milcGris}{Saber y saber hacer}
Al terminar podrás: (1) \textbf{identificar} los cinco errores más comunes de una tabla; (2) \textbf{aplicar} las herramientas de limpieza de Excel en el orden correcto; (3) \textbf{analizar} si la tabla limpia ya sirve para calcular; (4) \textbf{explicar} en una bitácora qué corregiste, cómo y por qué.
\end{softbox}



\small
\begin{tabularx}{\textwidth}{>{\bfseries}p{3.0cm}Y}
\rowcolor{milcGradePrimary!12}Autor & Dr. Álvaro Cárdenas Orozco\\
DBA & Tratamiento de datos cuantitativos con criterio ético (MEN, T\&I)\\
Referentes & Phronesis aristotélica · Floridi (ética del dato) · Estoicismo (ver lo que es)\\
Producto final & Una tabla de Excel con al menos 20 registros y 5 campos sobre un tema real del colegio, con formato de tabla aplicado. Y una bitácora que documente al menos cinco errores corregidos, con captura de antes y de después de cada uno.\\
\end{tabularx}
\normalsize

\puente{En la sesión 2 le pusiste tipo a cada columna. Hoy conviertes esas columnas en una tabla con campos y registros, y la limpias con bitácora. En la sesión 4 vas a calcular sobre ella, y solo sirve si hoy quedó limpia.}

\milcestacion{milcGradeDark}{Ruta MILC: así vamos a aprender}
\begin{tabularx}{\textwidth}{>{\centering\arraybackslash}X>{\centering\arraybackslash}X>{\centering\arraybackslash}X>{\centering\arraybackslash}X}
\phasecard{1}{milcVerde}{milcGris}{Escucha\\[-1mm]\small Observo, escucho y pregunto.} &
\phasecard{2}{milcTurquesa}{milcGris}{Entender\\[-1mm]\small Sistematizo y ordeno.} &
\phasecard{3}{milcMagenta}{milcGris}{Hacer\\[-1mm]\small Construyo y aplico.} &
\phasecard{4}{milcVino}{milcGris}{Cierre\\[-1mm]\small Evalúo y reflexiono.}
\end{tabularx}


\puente{Antes de limpiar, vas a mirar una tabla sucia que te entrega tu docente y a hacer el inventario de sus problemas. Todavía no corriges nada. Solo miras y anotas.}

\milcestacion{milcVerde}{Estación 1 de 3 · Escucha}

\begin{softbox}{milcVerde}{milcGris}{Escena para escuchar}
\textbf{Actividad 1 · IDENTIFICA --- Inventario de la suciedad} (15 min · individual). (1) Recibe la tabla con errores que te entrega o proyecta tu docente: una lista del grado con nombres, edades y fechas mal capturados. (2) Mira una columna a la vez. (3) Anota cada problema que veas: espacios de más, mayúsculas distintas, fechas raras, filas repetidas, celdas vacías. (4) Marca con una estrella los problemas que dañarían un cálculo. (5) Cuenta cuántos registros tiene la tabla y cuántos parecen sospechosos.
\end{softbox}

{\sffamily\bfseries\fontsize{8}{10}\selectfont\textcolor{milcNegro!55}{MARCA CUANDO LO HAYAS HECHO}}
\milccheckrow{Tengo una lista numerada de problemas, con la columna donde aparece cada uno.}
\milccheckrow{Marqué con estrella los que dañarían una suma o un promedio.}
\milccheckrow{Conté cuántos registros hay y cuántos son sospechosos.}

\begin{infoband}{milcVerde}


\textbf{Lo que va al cuaderno (Actividad 1 · IDENTIFICA).} \textbf{Título:} «Actividad 1 --- Inventario de la suciedad». \textbf{Formato:} lista numerada con problema, columna y marca de impacto (alto, medio, bajo). \textbf{Extensión:} media página. \textbf{Sabes que terminaste cuando} tienes al menos cinco problemas anotados y sabes por cuál empezarías.
\end{infoband}

\puente{Ya tienes el inventario. Ahora vas a entender los cinco errores comunes, la herramienta de Excel que arregla cada uno y el orden en que se aplican. Y vas a limpiar en pareja.}

\milcestacion{milcTurquesa}{Estación 2 de 3 · Entender}

\begin{softbox}{milcTurquesa}{milcGris}{Lo que vas a entender}
\textbf{Actividad 2 · APLICA --- Limpia la tabla en el orden correcto} (30 min · parejas). Un inventario sin herramientas es una lista de quejas. (1) Con tu pareja, abran la tabla sucia y guárdenla con otro nombre: la original no se toca. (2) Quiten los espacios de más con la función \texttt{=ESPACIOS()} o con Buscar y reemplazar. (3) Unifiquen las mayúsculas con \texttt{=NOMPROPIO()}. (4) Revisen las filas repetidas antes de quitarlas con Datos, Quitar duplicados. (5) Apliquen formato de fecha a la columna de fechas. (6) Decidan qué hacer con cada celda vacía, una por una. (7) Tomen captura antes y después de cada paso.
\end{softbox}

\milcpaso{1}{milcTurquesa}{\textbf{Los cinco errores y su herramienta.} Espacios de más: \texttt{=ESPACIOS(celda)} o Buscar y reemplazar. Mayúsculas distintas: \texttt{=NOMPROPIO()}, \texttt{=MAYUSC()} o \texttt{=MINUSC()}. Fechas mal escritas: formato de celda Fecha, como en la sesión 2. Filas repetidas: Datos, Quitar duplicados. Celdas vacías: se filtran y se decide caso por caso.}
\milcpaso{2}{milcTurquesa}{\textbf{El orden importa.} Primero los espacios, porque afectan a todo lo demás. Después las mayúsculas, porque sin unificarlas los duplicados no se detectan bien. Después los duplicados. Después las fechas. Al final las celdas vacías, una por una. En otro orden haces el trabajo dos veces.}
\milcpaso{3}{milcTurquesa}{\textbf{Una celda vacía no es un cero.} Si la llenas con cero, el promedio baja. Si la dejas vacía, Excel la ignora al promediar. Cada vacío se decide: rellenar con el dato real si lo sabes, dejar vacío si no, o quitar la fila si no sirve. Y se anota la decisión.}
\milcpaso{4}{milcTurquesa}{\textbf{Una tabla limpia no es la que se ve bonita: es la que sirve.} Sirve cuando las fórmulas funcionan, los filtros encuentran todo y los gráficos no engañan. Mide la limpieza por lo que vas a hacer después, no por cómo se ve ahora.}

\begin{drawbox}{milcTurquesa}{La tabla limpia, en una mirada}
\textbf{Sin espacios de más:} ningún carácter invisible. \\[1mm] \textbf{Mayúsculas iguales:} un criterio por columna. \\[1mm] \textbf{Fechas con formato:} se ordenan y se restan. \\[1mm] \textbf{Sin duplicados:} cada registro una sola vez. \\[1mm] \textbf{Vacíos decididos:} cada uno con su razón. \\[1mm] \textbf{Bitácora completa:} cuatro columnas por corrección.
\end{drawbox}

\begin{softbox}{milcMostaza}{milcGris}{Errores comunes y su corrección}
(1) \textbf{Limpiar sin anotar}: nadie puede verificar qué hiciste. (2) \textbf{Quitar duplicados sin revisar}: dos filas que parecen iguales pueden tener una diferencia importante. (3) \textbf{Llenar los vacíos con ceros}: baja los promedios. (4) \textbf{Limpiar en cualquier orden}: unas correcciones dependen de otras. (5) \textbf{No tomar la captura del antes}: pierdes la prueba del cambio.
\end{softbox}

\begin{infoband}{milcTurquesa}


\textbf{Lo que va al cuaderno (Actividad 2 · APLICA).} \textbf{Título:} «Actividad 2 --- La limpieza en orden». \textbf{Formato:} lista de los cinco pasos en el orden que los hicieron, con la herramienta usada al lado. \textbf{Extensión:} media página. \textbf{Sabes que terminaste cuando} los cinco pasos están en orden, cada uno con su herramienta, y tienen una captura de antes y de después por paso.
\end{infoband}

\puente{Ya limpiaste con tu pareja. Ahora escribes tú la bitácora de tu tabla y compruebas, con una fórmula de prueba, que la tabla ya sirve.}

\milcestacion{milcMagenta}{Estación 3 de 3 · Hacer}

\begin{softbox}{milcMagenta}{milcGris}{Cómo vas a construir tu producto}
\textbf{Actividad 3 · EXPLICA --- Escribe la bitácora y prueba la tabla} (25 min · individual). Ahora cuentas lo que hiciste, para que otro pueda repetirlo. (1) Abre tu tabla limpia y selecciónala completa. (2) Aplica Insertar, Tabla, para que Excel la reconozca como tabla. (3) Escribe la bitácora: por cada corrección, qué error había, qué herramienta usaste, cuándo y por qué. (4) Prueba la tabla con \texttt{=PROMEDIO()} sobre una columna numérica: si da error, quedó un número como texto. (5) Cuenta los registros finales y compáralos con los iniciales.
\end{softbox}

\milcpaso{1}{milcMagenta}{\textbf{La bitácora tiene cuatro columnas.} Error encontrado: «María aparece tres veces escrita distinto». Herramienta: «NOMPROPIO para unificar». Cuándo: fecha y hora. Justificación: «la mayoría ya estaba así». Sin la cuarta columna, nadie sabe por qué decidiste eso.}
\milcpaso{2}{milcMagenta}{\textbf{Escribir para alguien que no estuvo.} Tu bitácora la va a leer alguien que no vio tu pantalla. Cada línea debe bastar para repetir el paso. Prueba: dásela a tu pareja y pídele que repita una corrección solo leyendo.}
\milcpaso{3}{milcMagenta}{\textbf{La prueba de la fórmula.} Una tabla limpia acepta \texttt{=PROMEDIO()} sin error y devuelve un valor razonable. Si el promedio sale más bajo de lo esperado, hay ceros donde debía haber vacíos. Si sale error, hay texto donde debía haber número.}
\milcpaso{4}{milcMagenta}{\textbf{Cuenta lo que quitaste.} Registros iniciales menos registros finales es lo que eliminaste. Ese número va en la bitácora: si quitaste cinco filas y no sabes por qué, la tabla no está limpia; está recortada.}

\begin{drawbox}{milcMagenta}{Antes de entregar}
\checkbox Tabla original guardada aparte, sin tocar. \\[1mm] \checkbox Espacios, mayúsculas, duplicados, fechas y vacíos corregidos en ese orden. \\[1mm] \checkbox Formato de tabla aplicado (Insertar, Tabla). \\[1mm] \checkbox Bitácora con al menos cinco correcciones y sus cuatro columnas. \\[1mm] \checkbox Captura de antes y de después por cada corrección. \\[1mm] \checkbox \texttt{=PROMEDIO()} funciona sobre una columna numérica. \\[1mm] \checkbox Registros iniciales y finales anotados.
\end{drawbox}

\begin{softbox}{milcMagenta}{milcGris}{Plantilla de trabajo}
\textbf{Error.} \textit{Qué había mal, en una línea.} \\[1mm] \textbf{Herramienta.} \textit{Función, menú o atajo que usaste.} \\[1mm] \textbf{Cuándo.} \textit{Fecha y hora.} \\[1mm] \textbf{Justificación.} \textit{Por qué esa corrección y no otra.} \\[1mm] \textbf{Resultado.} \textit{Captura del después.}
\end{softbox}

\begin{infoband}{milcMagenta}


\textbf{Lo que va al cuaderno (Actividad 3 · EXPLICA).} \textbf{Título:} «Actividad 3 --- Mi bitácora de limpieza». \textbf{Formato:} tabla de 4 columnas (error / herramienta / cuándo / justificación) con al menos 5 filas, más la línea «registros iniciales y finales». \textbf{Extensión:} una página. \textbf{Sabes que terminaste cuando} tu pareja repite una corrección leyendo solo tu bitácora, y \texttt{=PROMEDIO()} devuelve un valor sin error.
\end{infoband}

\puente{Lo que entregas hoy: la tabla limpia con formato de tabla y la bitácora con cinco correcciones documentadas. La prueba de calidad: otra persona puede leer tu bitácora y repetir tu limpieza sin preguntarte nada.}

\milcestacion{milcMostaza}{Producto de la sesión}
\textbf{Tabla limpia con formato de tabla + bitácora de cinco correcciones + capturas}.
\begin{softbox}{milcMostaza}{milcGris}{Criterios de calidad}
(1) La tabla tiene al menos 20 registros y 5 campos, con formato de tabla. (2) Los cinco tipos de error se corrigieron en el orden indicado. (3) La bitácora tiene al menos cinco correcciones con sus cuatro columnas. (4) Hay captura de antes y de después por cada corrección. (5) Una fórmula de promedio funciona sobre la tabla sin error. (6) La bitácora dice cuántos registros había y cuántos quedaron.
\end{softbox}

\puente{Antes de cerrar, mira lo que hiciste desde cinco lados. Limpiar sin anotar es esconder el trabajo; anotar lo que corregiste es dejar que otros vuelvan a mirarlo.}

\milcestacion{milcVino}{Evaluación liberadora · cinco dimensiones}

\begin{softbox}{milcVino}{milcGris}{Sentido de la evaluación}
Evaluar no es solo revisar una nota. Es reconocer qué aprendí, qué cambió en mí, qué emociones manejé, cómo me relaciono con la ciudad y la comunidad, y qué le dejaré a quien viene detrás.
\end{softbox}

\textbf{Mi reflexión liberadora en cinco dimensiones.} Completa cada frase iniciada:

\small
\begin{tabularx}{\textwidth}{>{\bfseries\RaggedRight\arraybackslash}p{3.4cm}Y>{\RaggedRight\arraybackslash}p{4.6cm}}
\rowcolor{milcVino}\color{white}Dimensión & \color{white}\bfseries Mi respuesta (frame starter) & \color{white}\bfseries Algo que descubrí\\
1. Personal & Lo que más cambió en mí fue \linead{6.5cm} & \linead{4.2cm}\\
2. Emocional & La emoción más fuerte fue \linead{2.5cm} y la regulé \linead{3cm} & \linead{4.2cm}\\
3. Ciudadana & En democracia esto importa porque \linead{6cm} & \linead{4.2cm}\\
4. Local (Cartago) & En mi barrio sería útil para \linead{6.5cm} & \linead{4.2cm}\\
5. Intergeneracional & A mi abuela/abuelo le diría \linead{6.5cm} & \linead{4.2cm}\\
\end{tabularx}
\normalsize

\begin{infoband}{milcVino}
\textbf{Para la próxima sustentación.} Vuelve a esta tabla en seis meses. Verás que las respuestas cambian — no porque lo aprendido cambie, sino porque tú cambias. La evaluación liberadora es una conversación contigo a través del tiempo.
\end{infoband}


\puente{Para terminar, tres ideas sobre volver a mirar, contadas con palabras de esta guía: Dussel sobre creer que la información se entiende sola, Marco Aurelio sobre aceptar la corrección, Floridi sobre saber qué datos guardar y cuáles no.}

\milcestacion{milcGradeDark}{Tres ideas para cerrar · Dussel, un estoico y Floridi}

\begin{softbox}{milcVino}{milcGris}{Enrique Dussel · lente del nosotros}
\textit{Es ingenuo creer que la información se entiende sola, como si no hubiera conflictos detrás de cada dato.}

{\footnotesize\itshape\color{milcNegro!70}--- idea tomada de Enrique Dussel, contada con palabras de esta guía. No es una cita textual.}

Una tabla «limpia» esconde decisiones: qué fila quitaste, qué vacío llenaste, a quién dejaste fuera. La bitácora es la forma de no esconderlas. Sin ella, tu tabla parece neutral y no lo es.

\textbf{¿Qué decisión de mi limpieza cambiaría el resultado si otra persona la hubiera tomado distinto?}
\end{softbox}
\linead{\linewidth}\\[1mm]
\linead{\linewidth}

\begin{softbox}{milcMostaza}{milcGris}{Estoicismo · Marco Aurelio · cuidado interior}
\textit{Cambiar de parecer y seguir el aviso de quien te corrige es tan libre como mantenerte en lo tuyo.}

{\footnotesize\itshape\color{milcNegro!70}--- idea tomada de Marco Aurelio, contada con palabras de esta guía. No es una cita textual.}

Cuando tu pareja lee tu bitácora y dice «este paso no se entiende», tienes dos opciones: defenderte o corregir. La segunda es la que hace mejor la tabla. Y la bitácora existe precisamente para que alguien pueda corregirte.

\textbf{¿Escribí la bitácora para que me corrijan, o para que parezca que no me equivoqué?}
\end{softbox}
\linead{\linewidth}\\[1mm]
\linead{\linewidth}

\begin{softbox}{milcTurquesa}{milcGris}{Luciano Floridi · lente de la infoesfera}
\textit{La mitad de los datos que guardamos es basura; el problema es que no sabemos cuál mitad.}

{\footnotesize\itshape\color{milcNegro!70}--- idea tomada de Luciano Floridi, contada con palabras de esta guía. No es una cita textual.}

Limpiar una tabla es decidir cuál mitad sirve. Cada fila repetida, cada vacío, cada «MARIA» con mayúsculas es una decisión sobre qué cuenta y qué no. Por eso se anota.

\textbf{De todo lo que quité o cambié, ¿qué parte podría haber sido un dato real que perdí?}
\end{softbox}
\linead{\linewidth}\\[1mm]
\linead{\linewidth}

\begin{infoband}{milcNegro}
\textbf{Nota para el docente.} Las tres ideas están contadas con palabras de esta guía a partir del pensamiento de cada autor; no son citas textuales. De 9.º en adelante se trabajan con la obra y su referencia.
\end{infoband}

\milcestacion{milcVino}{Mi autoevaluación · marca una casilla por fila}

{\small
\setlength{\extrarowheight}{2.2pt}
\begin{tabularx}{\textwidth}{>{\RaggedRight\arraybackslash\bfseries}p{3.0cm}YYY}
\rowcolor{milcVino}\color{white}\bfseries Criterio & \color{white}\bfseries Lo logré & \color{white}\bfseries En proceso & \color{white}\bfseries Necesito apoyo\\[.6mm]
Comprensión del tema & \milcopcion{Explico las ideas con mis palabras.} & \milcopcion{Reconozco las ideas, falta precisión.} & \milcopcion{Necesito repasar lo básico.}\\[.8mm]
\rowcolor{milcGris}Pensamiento computacional & \milcopcion{Usé los pilares al armar mi producto.} & \milcopcion{Usé algunos pilares.} & \milcopcion{Necesito releer la estación 2.}\\[.8mm]
Aplicación y producto & \milcopcion{Mi producto cumple los criterios.} & \milcopcion{Está armado, con ajustes pendientes.} & \milcopcion{Aún está en construcción.}\\[.8mm]
\rowcolor{milcGris}Reflexión 5 dimensiones & \milcopcion{Respondí cada una con honestidad.} & \milcopcion{Respondí algunas con detalle.} & \milcopcion{Mis respuestas son muy generales.}\\[.8mm]
Triángulo de pensamiento & \milcopcion{Conecté las 3 voces con mi trabajo.} & \milcopcion{Conecté al menos una.} & \milcopcion{No las relaciono con mi proceso.}\\[.8mm]
\end{tabularx}}

\vspace{3mm}
\textbf{Hoy entendí que:}\\[1mm]
\linead{\linewidth}\\[1mm]
\linead{\linewidth}

\textbf{Mi compromiso para la próxima sesión es:}\\[1mm]
\linead{\linewidth}\\[1mm]
\linead{\linewidth}

\begin{softbox}{milcMostaza}{milcGris}{Cita de despedida · Marco Aurelio}
\textit{``El obstáculo en el camino se vuelve el camino. Lo que estorba al camino se vuelve camino.''}\\[1mm]
Cada error de hoy, bien observado, se convierte en pieza del camino que pisarás mañana.
\end{softbox}

\small
\begin{tabularx}{\textwidth}{>{\bfseries}p{3.6cm}Y}
\rowcolor{milcGradeDark!12}Nombre completo: & \linead{10cm}\\
Fecha: & \linead{4cm} \quad Curso/grupo: \linead{4cm}\\
Mi firma: & \linead{9cm}\\
\end{tabularx}
\normalsize

\begin{infoband}{milcGradeDark}
\textbf{Para la próxima sesión.} Llega con tu tabla limpia y la bitácora completa. En la sesión 4 vas a calcular SUMA, PROMEDIO, MAX y MIN sobre esa tabla y a tomar una decisión con cada cifra. Si hoy quedó sucia, mañana esas cifras engañan.
\end{infoband}

\end{document}
